problem:  
Let \( M^{4n} \) be a closed, simply connected Riemannian manifold with **positive sectional curvature**, admitting a smooth **isometric \(S^1\)-action** with **nonempty fixed-point set**.  
Assume that every connected component of \(M^{S^1}\) is a **totally geodesic**, even‑codimensional submanifold, all of **codimension exactly \(2m \ge 8\)**.  
Let \(F\) be one such fixed-point component, and let \(\nu_F\) denote its normal bundle.  

By Wilking’s connectedness lemma and Kennard’s periodicity theorem, we know that if these submanifolds occur in **codimension ≤ 6**, then \(M\) has **4‑periodic rational cohomology**. No result is presently known in higher codimensions.

**Question:**  
If \(M\) admits such an isometric circle action with totally geodesic fixed‑point components of *minimum codimension \(2m = 8\)*, does \(M\) necessarily have 8‑periodic rational cohomology?  
Equivalently, does the existence of an isometric \(S^1\)-action with all fixed components of codimension 8 in a positively curved \(4n\)-manifold still force the rational cohomology ring to be generated by a degree‑8 element?  

If not, can one construct or rule out a counterexample where positive curvature coexists with failure of 8‑periodicity despite such symmetry and codimension control?

---

Why is it a "good" problem:  
This problem isolates the **next unknown threshold** in the interaction of positive curvature, isometric group actions, and cohomological rigidity.  
The classical machinery of Wilking and Kennard explains full 4‑periodicity when the circle‑fixed sets have codimension ≤ 6, but those arguments break precisely at codimension 8, where topological periodicity can no longer be derived from existing connectedness estimates.  

By insisting that *all* fixed-point components share codimension 8 and are totally geodesic, this question probes the first regime **beyond the reach of known theorems**. Resolving it would demand new geometric ideas, possibly linking **curvature‑induced connectivity**, **equivariant Morse theory**, and **rational homotopy periodicity** in ways not yet understood.  
The statement is conceptually simple—just positive curvature, a circle symmetry, and all fixed sets of the same codimension—but its resolution would represent a genuine advance on one of the central open frontiers surrounding the **Hopf and symmetry–rigidity conjectures**, making it a concise yet profound research-level problem.